\begin{env}[3.8.2]
\label{0.3.8.2-fr}
Sean $\mathcal{F}$, $\mathcal{G}$ dos haces de conjuntos (resp. de grupos,
de anillos) sobre $X$, y $u:\mathcal{F}\to\mathcal{G}$ un homomorfismo.
Entonces $u$ es también un homomorfismo \emph{continuo}
$\mathcal{F}'\to\mathcal{G}'$, donde $\mathcal{F}'$ y $\mathcal{G}'$ designan
los haces pseudodiscretos asociados a $\mathcal{F}$ y $\mathcal{G}$.
Esto se deduce, en efecto, de (3.2.5).
\end{env}
